\documentclass{article}

% Language setting
% Replace `english' with e.g. `spanish' to change the document language
%\usepackage[english,portuguese]{babel}
\usepackage[brazil]{babel}
\usepackage[utf8]{inputenc}

% Set page size and margins
% Replace `letterpaper' with `a4paper' for UK/EU standard size
\usepackage[a4paper,top=2cm,bottom=2cm,left=4cm,right=4cm,marginparwidth=1.75cm]{geometry}

% Useful packages
\usepackage{amsmath,float}
\usepackage{graphicx}
\usepackage[colorlinks=true, allcolors=blue]{hyperref}
\usepackage{diagbox}

%\newenvironment{namebox}
%{\begin{flushleft}Nome: \underline{\hspace{6cm}}\end{flushleft}} % Adjust the length of the underline as needed


\date{--/--/----} % Removendo a data

\begin{document}

\begin{flushleft}
	Universidade Federal da Bahia\\
	Instituto de Matemática e Estatística\\
	Departamento de Estatística\\
	MATF59 T01 - Estatística básica para humanidades - 2025.1\\
	Professor: Rodney Fonseca \\
%	Data: 24/09/2024
\end{flushleft}

\begin{center}
	\begin{Large}
		Fórmulas \\
	\end{Large}	
\end{center}


\bigskip


\begin{Large}
\noindent
Se o número de casos de uma categoria é $f$ e o número total de casos é $n$, então
\begin{align*}
\text{Proporção da categoria} = \frac{f}{n}
\end{align*}
e
\begin{align*}
\text{Porcentagem da categoria} = \left( 100 \times\frac{f}{n} \right) \%
\end{align*}
\end{Large}

\bigskip
--------------------------------------------


\begin{Large}
\noindent
Se o número de casos de uma categoria é $f_1$ e o número de casos numa outra categoria é $f_2$, então
\begin{align*}
\text{Razão da categoria 1 para a categoria 2} = \frac{f_1}{f_2}.
\end{align*}
Usando um número $n$ como base, dizemos que há $n \times \frac{f_1}{f_2} $ elementos da categoria 1 para cada $n$ elementos da categoria~2.
\end{Large}

\bigskip
--------------------------------------------


\begin{Large}
\noindent
Se o valor da variável no tempo 1 é $f_1$ e o valor no tempo 2 é $f_2$, então
\begin{align*}
\text{Taxa de variação do tempo 1 para o tempo 2} = 100 \times \left(\frac{f_2 - f_1}{f_1}\right)\%
\end{align*}
\end{Large}


\bigskip
--------------------------------------------

\begin{align*}
\text{Mediana} = \left\{\begin{array}{lc}
\text{valor na posição $\frac{n+1}{2}$}, & \text{se $n$ for ímpar} \\
\text{média dos valores nas posições $\frac{n}{2}$ e $\frac{n}{2}+1$}, & \text{se $n$ for par}
\end{array}\right.
\end{align*}

\begin{Huge}
\begin{center}
Se temos $n$ valores $x_1, x_2, \ldots, x_n$, a média deles será
\begin{align*}
\bar{x} = \frac{x_1 + x_2 + \cdots + x_n}{n} = \frac{\sum_{i=1}^n x_i}{n}
\end{align*}
\end{center}
\end{Huge}

\bigskip


\begin{Large}
\begin{center}
Sendo $x_{\max}$ e $x_{\min}$ o máximo e mínimo dos dados, respectivamente, a amplitude total será
\begin{align*}
AT = x_{\max} - x_{\min}
\end{align*}
\end{center}
\end{Large}


\bigskip


\begin{Large}
\begin{center}
Sejam $x_1, \ldots, x_n$ os $n$ valores e $\bar{x}$ a média deles. \\ A variância será:
\begin{align*}
var(X) = \frac{\sum_{i=1}^n (x_i - \bar{x})^2}{n} = \frac{(x_1 - \bar{x})^2 + (x_1 - \bar{x})^2 + \cdots + (x_n - \bar{x})^2}{n}
\end{align*}
\end{center}
\end{Large}


\bigskip


\begin{Large}
\begin{center}
O desvio padrão dos $n$ valores $x_1, \ldots, x_n$ será:
\begin{align*}
dp(X) = \sqrt{ \frac{\sum_{i=1}^n (x_i - \bar{x})^2}{n} } = \sqrt{ \frac{(x_1 - \bar{x})^2 + (x_1 - \bar{x})^2 + \cdots + (x_n - \bar{x})^2}{n} }
\end{align*}
\end{center}
\end{Large}


\bigskip


\begin{Large}
\begin{center}
A medida de assimetria de $x_1, \ldots, x_n$ é:
\begin{align*}
\frac{\frac{1}{n}\sum_{i=1}^n (x_i - \bar{x})^3}{var(X)^{3/2}} = \frac{\frac{1}{n}\sum_{i=1}^n (x_i - \bar{x})^3}{\left( \frac{1}{n}\sum_{i=1}^n (x_i - \bar{x})^2 \right)^{3/2}},
\end{align*}
em que $\bar{x}$ e $var(X)$ são a média variância amostrais
\end{center}
\end{Large}


\bigskip


\begin{Large}
\begin{center}
A medida de curtose de $x_1, \ldots, x_n$ é:
\begin{align*}
\frac{\frac{1}{n}\sum_{i=1}^n (x_i - \bar{x})^4}{var(X)^{2}} = \frac{\frac{1}{n}\sum_{i=1}^n (x_i - \bar{x})^4}{\left( \frac{1}{n}\sum_{i=1}^n (x_i - \bar{x})^2 \right)^{2}},
\end{align*}
em que $\bar{x}$ e $var(X)$ são a média variância amostrais
\end{center}
\end{Large}

\bigskip


\begin{Large}
\begin{center}
Sejam $\bar{x}$ e $dp(X)$ a média e o desvio padrão dos $n$ valores $x_1, \ldots, x_n$. O coeficiente de variação é:
\begin{align*}
CV(X) = \frac{dp(X)}{\bar{x}}
\end{align*}
\end{center}
\end{Large}


\bigskip


\begin{minipage}{0.8\textwidth}
Tabela: Distribuição das frequências absolutas segundo sexo e curso
\begin{center}
\begin{tabular}{c|cc}
\hline
\diagbox{Curso}{Sexo} & Masculino & Feminino \\
\hline
Física    & 100 & 20 \\
Geografia & 40  & 40 \\
\hline
\end{tabular}
\end{center}
\end{minipage}


\bigskip


\begin{minipage}{0.7\textwidth}
Tabela: Distribuição das frequências absolutas segundo sexo e curso, com frequências relativas em relação ao curso
\begin{tabular}{c|cc|c}
\hline
\diagbox{Curso}{Sexo} & Masculino & Feminino & Total\\
\hline
Física    & 100 (71\%) & 20 (33\%) & 120 (60\%) \\
Geografia & 40 (29\%)  & 40 (67\%) &  80 (40\%) \\
\hline
Total & 140 (100\%) & 60 (100\%) &  200 (100\%) \\
\hline
\end{tabular}
\end{minipage}


\bigskip


\begin{minipage}{0.7\textwidth}
\begin{center}
\begin{tabular}{c|cc|c}
\hline
 & Recuperados & Mortes & Total\\
\hline
Vacinados     & 197 &  2 & 199 \\
Não-vacinados & 139 & 19 & 158 \\
\hline
Total & 336 & 21 & 357 \\
\hline
\end{tabular}\\
{\tiny Fonte: Yule, G. U. (1912). On the methods of measuring association between two attributes. JRSS, 75(6), 579-652.}
\end{center}
\end{minipage}





\end{document}